\documentclass[11pt,a4paper]{article}

\usepackage[margin=1.8cm,top=1.6cm,bottom=1.6cm]{geometry}
\usepackage[T1]{fontenc}
\usepackage{lmodern}
\usepackage{enumitem}
\usepackage{array}
\usepackage{xcolor}
\usepackage{fancyhdr}

\newcommand{\cmd}[1]{\texttt{#1}}

\pagestyle{fancy}
\fancyhf{}
\fancyhead[L]{Lab Exam: Git \& GitHub}
\fancyhead[R]{MSIS DevOps}
\fancyfoot[C]{\thepage}
\renewcommand{\headrulewidth}{0.4pt}

\begin{document}

\begin{center}
  {\LARGE\bfseries Lab Exam: Git \& GitHub Workflow}\\[6pt]
  {\large MSIS DevOps}
\end{center}

\vspace{4pt}

\noindent
\begin{tabular}{@{}p{0.5\textwidth}p{0.45\textwidth}@{}}
  \textbf{Duration:} 90 minutes (practical) & \textbf{Followed by:} Viva voce \\[4pt]
  \textbf{Date:} 13 October 2026 & \\
\end{tabular}

\vspace{8pt}
\hrule
\vspace{8pt}

\subsection*{Instructions}
\begin{itemize}[leftmargin=*,itemsep=2pt]
  \item Perform all operations in order. Each step builds on the previous one.
  \item Use the Git command line for all local operations.
  \item Keep your terminal open. You may be asked to show your command history,
        \cmd{git log}, and the pull requests on GitHub during the viva.
  \item Make one commit per step that asks for a commit, with a meaningful commit message.
  \item The practical is followed by a viva on the commands you used and why.
\end{itemize}

\subsection*{Tasks}
Perform the following Git operations:

\begin{enumerate}[label=\textbf{\arabic*.},leftmargin=*,itemsep=3pt]
  \item Create a local repository called \cmd{msis-devops} containing a
        \cmd{readme.md} file, and push it to a \textbf{private} GitHub repository
        with the same name.

  \item Create a new branch called \cmd{project-setup}. Set up a sample Python
        project containing \cmd{demo.py} and an appropriate \cmd{.gitignore} file.

  \item Commit the changes, push the branch to the remote repository, and create
        a pull request (PR).

  \item Merge the PR into \cmd{main}.

  \item Switch to \cmd{main} and pull the latest changes from the remote
        repository. Create a new branch called \cmd{math-ops}. In \cmd{demo.py},
        add functions named \cmd{add} and \cmd{sub}, each accepting two integers
        and returning the result.

  \item Commit the changes, push the branch, create a PR, and merge it into
        \cmd{main}.

  \item Switch to \cmd{main} and pull the latest changes. Create a new branch
        called \cmd{multiply-op} and implement a \cmd{multiply} function in
        \cmd{demo.py}. Commit the changes, push the branch, and create a PR.
        \textbf{Do not merge this PR yet.}

  \item Switch back to \cmd{main} and create a new branch called
        \cmd{divide-op}. Implement a \cmd{divide} function in \cmd{demo.py}.
        Commit the changes, push the branch, and create a PR.

  \item Merge the PR created in step 7 into \cmd{main}.

  \item Switch to \cmd{divide-op} and pull the latest changes from the remote
        \cmd{main} branch into it.

  \item Resolve any merge conflicts, ensuring that both the \cmd{multiply} and
        \cmd{divide} functions are retained in \cmd{demo.py}.

  \item Commit the conflict resolution and push the updated \cmd{divide-op}
        branch to the remote repository.

  \item Merge the PR created in step 8 into \cmd{main}.

  \item Switch to \cmd{main}, pull the latest changes, and display \cmd{demo.py}
        to confirm that it contains all four functions: \cmd{add}, \cmd{sub},
        \cmd{multiply}, and \cmd{divide}.
\end{enumerate}

\end{document}
